%% kapa: exam=1 round=37 year=2026 period=2 form=A subject=appraisal_law subject_ko=감정평가관계법규 seq=39 number=39
%% source: Q-Net 감정평가사 기출문제 (cst00309) artlSeq=5258924 file=raw/1008/5258924/2026년 제37회 감정평가사 1차 2교시.pdf
%% conversion: kordoc 4.24.51 -> tools/gen1.py (auto); review=proofread-sonnet; flags=none
\begin{question}{39}
도시 및 주거환경정비법령상 조합설립추진위원회의 구성ㆍ승인에 관한 조문의 일부이다. ( )에 들어갈 내용으로 옳은 것은?
\begin{framedtext}
조합을 설립하려는 경우에는 다음 각 호의 사항에 대하여 토지등소유자 ( ㄱ )의 동의를 받아 조합설립을 위한 추진위원회를 구성하여 국토교통부령으로 정하는 방법과 절차에 따라 시장ㆍ군수등의 승인을 받아야 한다. \par
\hspace*{1em}1. 추진위원회 위원장을 포함한 ( ㄴ )명 이상의 추진위원회 위원 \par
\hspace*{1em}2. < 생략 > \par
\end{framedtext}
\begin{choices}
  \item ㄱ: 과반수, ㄴ: 5
  \item ㄱ: 과반수, ㄴ: 7
  \item ㄱ: 3분의 2 이상, ㄴ: 5
  \item ㄱ: 3분의 2 이상, ㄴ: 7
  \item ㄱ: 3분의 2 이상, ㄴ: 10
\end{choices}
\end{question}
\answer{1}
